\begingroup
\clearpage\SourceLeaf{341}\ArithmeticHead{334}{L}{ANALYSIS.\quad NACHLASS.}
\NumberedItem{4}
\noindent{\small[S. 244]}\par
\[\begin{aligned}
\Pi\,0{,}2&=x\\
\Pi\,0{,}4&=y\\
\Pi\,0{,}6&=0{,}24\pi:y\sin72^0\\
\Pi\,0{,}8&=0{,}16\pi:x\sin36^0\\
\Pi\,9&=10{.}189{.}192\\
\Pi\,9{,}2&=x{.}1066{.}1023{.}966{.}768:10^9\\
\Pi\,9{,}4&=y{.}192^3{.}329{.}561{.}777\\
\Pi\,9{,}6&=\pi{.}96^4{.}462{.}988{.}989:y\sin72^0\\
\Pi\,9{,}8&=\pi{.}672^3{.}528{.}493{.}494:x\sin36^0\\
\Pi\,10{,}0&=100{.}189{.}192.
\end{aligned}\]
\par\indent [Die Logarithmen werden auf 12 Dezimalstellen gerechnet mit Benutzung der Werte:]
\[\begin{aligned}
[\log_{10}]\,x\;&{.}\;{.}\;{.}\;9{,}962922504818,\\
{}[\log_{10}]\,y\;&{.}\;{.}\;{.}\;9{,}948052773504,
\end{aligned}\]
\noindent [die aber (vergl. Werke III, S. 161) von der 9. Ziffer an unrichtig sind.]
\par\indent Invenimus autem
\[\begin{aligned}
F(-\tfrac15,-\tfrac15,10,1)&=1{,}0041288636
=\frac{\Pi9{.}\Pi9{,}4}{\Pi9{,}2{.}\Pi9{,}2}\\[.8em]
F(-\tfrac15,-\tfrac15,10\tfrac15,1)&=1{,}0040454154
=\frac{\Pi9{,}2{.}\Pi9{,}6}{\Pi9{,}4{.}\Pi9{,}4}\\[.8em]
F(\tfrac15,\tfrac15,11,1)&=1{,}0038882487
=\left[\frac{\Pi10{.}\Pi9{,}6}{\Pi9{,}8{.}\Pi9{,}8}\right].
\end{aligned}\]
\begin{center}\rule{.22\linewidth}{.4pt}\end{center}
\[\begin{aligned}
\Lambda z&=\log1+\log2+\log3+{.}{.}{.}+\log k+z\log k\\
&\quad-\log(z+1)-\log(z+2)-{.}{.}{.}-\log(z+k)\\[.8em]
\Lambda'z&=-\frac1{z+1}-\frac1{z+2}-\frac1{z+3}-{.}{.}{.}-\frac1{z+k}+\log k\\[.8em]
\Lambda''z&=\frac1{(z+1)^2}+\frac1{(z+2)^2}+\frac1{(z+3)^2}+\text{etc.}\\[.8em]
\Lambda'''z&=-\frac2{(z+1)^3}-\frac2{(z+2)^3}-\frac2{(z+3)^3}-{.}{.}{.}
\end{aligned}\]
\begin{center}etc.\end{center}
\clearpage\SourceLeaf{342}\ArithmeticHead{335}{R}{{\scriptsize ZUR THEORIE DER REIHE $F(\alpha,\beta,\gamma,x)$.\quad II.}}
\[\begin{aligned}
\Lambda(z+\omega)&=\Lambda z+\omega\left(1+\frac12+\frac13+\text{etc.}+\frac1z-0{,}[57721\;{.}\;{.}\;{.}]\right)\\[.8em]
&\quad+\frac12\omega\omega\left(\frac1{(z+1)^2}+\frac1{(z+2)^2}+\frac1{(z+3)^2}+\text{etc.}\right)\\[.8em]
&\quad-\frac13\omega^3\left(\frac1{(z+1)^3}+\frac1{(z+2)^3}+\frac1{(z+3)^3}+\text{etc.}\right)\\[.8em]
&\quad+\frac14\omega^4\left(\frac1{z+1)^4}+\frac1{(z+2)^4}+\frac1{(z+3)^4}+\text{etc.}\right)
\end{aligned}\]
\[ {.}\quad{.}\quad{.}\quad{.}\quad{.}\quad{.}\quad{.}\quad{.}\quad{.}\quad{.}\quad{.}\quad{.}\quad{.} \]
\noindent [Auf S. 245 folgt eine Tafel der Werte von $\log_{10}\Pi z$ für die Werte von $z$ von $9{,}20$ bis zu $10{,}20$, fortschreitend um je $0{,}02$; die Tafel ist auf 10 Dezimalstellen gerechnet, aber mit Benutzung der ungenauen Werte von $\log x$ und $\log y$.]
\clearpage
\endgroup
